\subsection{Deficit por radio para \texttt{HIP 97166 c} en \texttt{cube12\_cluster0}}
\begin{table}[H]
\centering
\small
\caption{Conteos observados, esperados y deficit local para el planeta ancla \texttt{HIP 97166 c}.}
\begin{tabular}{lrrrr}
\toprule
Radio & $N_{\mathrm{obs}}$ & $N_{\mathrm{exp}}$ & $\Delta N$ & $\Delta_{\mathrm{rel}}$ \\
\midrule
$r_{\mathrm{kNN}}$ & 10.00 & 11.00 & 1.00 & 0.091 \\
$r_{\mathrm{node\_median}}$ & 1617.00 & 1144.00 & -473.00 & -0.413 \\
$r_{\mathrm{node\_q75}}$ & 2550.00 & 1914.00 & -636.00 & -0.332 \\
\bottomrule
\end{tabular}
\end{table}
El deficit depende de la escala local elegida. Esto sugiere que el ancla es util para priorizacion, pero la evidencia es sensible al radio y debe considerarse exploratoria.

\subsection{Deficit por radio para \texttt{HIP 90988 b} en \texttt{cube17\_cluster2}}
\begin{table}[H]
\centering
\small
\caption{Conteos observados, esperados y deficit local para el planeta ancla \texttt{HIP 90988 b}.}
\begin{tabular}{lrrrr}
\toprule
Radio & $N_{\mathrm{obs}}$ & $N_{\mathrm{exp}}$ & $\Delta N$ & $\Delta_{\mathrm{rel}}$ \\
\midrule
$r_{\mathrm{kNN}}$ & 10.00 & 11.00 & 1.00 & 0.091 \\
$r_{\mathrm{node\_median}}$ & 8.00 & 9.00 & 1.00 & 0.111 \\
$r_{\mathrm{node\_q75}}$ & 62.00 & 63.00 & 1.00 & 0.016 \\
\bottomrule
\end{tabular}
\end{table}
El deficit aparece en las tres escalas locales, por lo que la evidencia topologica de submuestreo es mas estable. Aun asi, no debe interpretarse como conteo absoluto de objetos reales ausentes.

\subsection{Deficit por radio para \texttt{HD 42012 b} en \texttt{cube19\_cluster5}}
\begin{table}[H]
\centering
\small
\caption{Conteos observados, esperados y deficit local para el planeta ancla \texttt{HD 42012 b}.}
\begin{tabular}{lrrrr}
\toprule
Radio & $N_{\mathrm{obs}}$ & $N_{\mathrm{exp}}$ & $\Delta N$ & $\Delta_{\mathrm{rel}}$ \\
\midrule
$r_{\mathrm{kNN}}$ & 10.00 & 11.00 & 1.00 & 0.091 \\
$r_{\mathrm{node\_median}}$ & 120.00 & 121.00 & 1.00 & 0.008 \\
$r_{\mathrm{node\_q75}}$ & 293.00 & 294.00 & 1.00 & 0.003 \\
\bottomrule
\end{tabular}
\end{table}
El deficit aparece en las tres escalas locales, por lo que la evidencia topologica de submuestreo es mas estable. Aun asi, no debe interpretarse como conteo absoluto de objetos reales ausentes.

\subsection{Deficit por radio para \texttt{HD 42012 b} en \texttt{cube26\_cluster6}}
\begin{table}[H]
\centering
\small
\caption{Conteos observados, esperados y deficit local para el planeta ancla \texttt{HD 42012 b}.}
\begin{tabular}{lrrrr}
\toprule
Radio & $N_{\mathrm{obs}}$ & $N_{\mathrm{exp}}$ & $\Delta N$ & $\Delta_{\mathrm{rel}}$ \\
\midrule
$r_{\mathrm{kNN}}$ & 10.00 & 11.00 & 1.00 & 0.091 \\
$r_{\mathrm{node\_median}}$ & 120.00 & 121.00 & 1.00 & 0.008 \\
$r_{\mathrm{node\_q75}}$ & 293.00 & 294.00 & 1.00 & 0.003 \\
\bottomrule
\end{tabular}
\end{table}
El deficit aparece en las tres escalas locales, por lo que la evidencia topologica de submuestreo es mas estable. Aun asi, no debe interpretarse como conteo absoluto de objetos reales ausentes.

\subsection{Deficit por radio para \texttt{HD 4313 b} en \texttt{cube17\_cluster10}}
\begin{table}[H]
\centering
\small
\caption{Conteos observados, esperados y deficit local para el planeta ancla \texttt{HD 4313 b}.}
\begin{tabular}{lrrrr}
\toprule
Radio & $N_{\mathrm{obs}}$ & $N_{\mathrm{exp}}$ & $\Delta N$ & $\Delta_{\mathrm{rel}}$ \\
\midrule
$r_{\mathrm{kNN}}$ & 10.00 & 11.00 & 1.00 & 0.091 \\
$r_{\mathrm{node\_median}}$ & 5.00 & 6.00 & 1.00 & 0.167 \\
$r_{\mathrm{node\_q75}}$ & 11.00 & 12.00 & 1.00 & 0.083 \\
\bottomrule
\end{tabular}
\end{table}
El deficit aparece en las tres escalas locales, por lo que la evidencia topologica de submuestreo es mas estable. Aun asi, no debe interpretarse como conteo absoluto de objetos reales ausentes.

